% Part: incompleteness
% Chapter: incompleteness-provability
% Section: fixed-point-lemma

\documentclass[../../../include/open-logic-section]{subfiles}

\begin{document}

\olfileid{inc}{inp}{fix}

\olsection{Lema Titik Tetep}

\begin{explain}
Lema titik tetep nyatakake yen kanggo saben !!{formula}~$!B(x)$, ana
!!a{sentence}~$!A$ saengga $\Th{T} \Proves !A \liff !B(\gn{!A})$,
yen $\Th{T}$ ngembangake~$\Th{Q}$. Ing kasus ukara paradoks tukang
ngapusi, kita kepengin $!A$ ekuivalen (bisa dibuktekake ing~$\Th{T}$)
karo~``$\gn{!A}$ salah,'' yaiku pratelan yen $\Gn{!A}$ iku
bilangan G\"odel saka !!{sentence} kang salah. Kanggo mangerteni
gagasan buktine, migunani yen dibandhingake karo andharan informal
Quine ngenani~$!A$ minangka
``{}`ngasilake pratelan salah yen didhisiki petikane' ngasilake
pratelan salah yen didhisiki petikane.'' Tumindak njupuk ekspresi
banjur mbentuk ukara kanthi ndhisiki ekspresi iku nganggo petikane
dhewe bisa diarani \emph{ndiagonalisasi} ekspresi iku, lan asile
diarani diagonalisasine. Dadi diagonalisasi saka
`ngasilake pratelan salah yen didhisiki petikane' yaiku
``{}`ngasilake pratelan salah yen didhisiki petikane' ngasilake
pratelan salah yen didhisiki petikane.'' Gatekna saiki yen ukara
paradoks tukang ngapusi saka Quine dudu diagonalisasi saka
`ngasilake pratelan salah' nanging saka
`ngasilake pratelan salah yen didhisiki petikane.' Dadi sipat kang
didiagonalisasi kanggo ngasilake ukara paradoks iku dhewe wis
ngemot diagonalisasi!{}

Ing basa aritmetika, kita mbentuk petikan saka !!a{formula} kang
nduweni siji variabel bebas kanthi ngitung bilangan G\"odele banjur
nglebokake numeral standar kanggo bilangan G\"odel mau ing variabel
bebas. Diagonalisasi saka~$!E(x)$ iku $!E(\num{n})$, ing ngendi
$n = \Gn{!E(x)}$. (Wiwit saiki, kita cekak $\num{\Gn{!E(x)}}$ dadi
$\gn{!E(x)}$.) Dadi yen $!B(x)$ iku ``iku salah,'' mula ``ngasilake
pratelan salah yen didhisiki petikane'' dadi ``ngasilake pratelan
salah yen ditrapake marang bilangan G\"odel saka diagonalisasine.''
Yen kita nduweni simbol~$\Obj{diag}$ kanggo fungsi $\fn{diag}(n)$ kang
ngitung bilangan G\"odel saka diagonalisasi !!{formula} kang nduweni
bilangan G\"odel~$n$, kita bisa nulis $!E(x)$ minangka
$!B(\Obj{diag}(x))$. Versi ukara paradoks saka Quine banjur dadi
diagonalisasine, yaiku $!E(\gn{!E(x)})$ utawa
$!B(\Obj{diag}(\gn{!B(\Obj{diag}(x))}))$. Mesthi wae, $!B(x)$ saiki
bisa dadi sipat liyane sembarang, lan konstruksi kang padha tetep
lumaku. Kanggo teorema ora jangkep, kita pilih $!B(x)$ minangka
``$x$~ora !!{derivable} ing~$\Th{T}$.'' Mula $!E(x)$ dadi
``ngasilake !!a{sentence} kang ora !!{derivable} ing~$\Th{T}$
yen ditrapake marang bilangan G\"odel saka diagonalisasine.''

Kanggo ngformalisasi iki ing~$\Th{T}$, kita kudu nemokake cara
ngformalisasi $\fn{diag}$. Fungsi $\fn{diag}(n)$ komputabel, malah
rekursif primitif: yen $n$ iku bilangan G\"odel saka
formula~$!E(x)$, $\fn{diag}(n)$ ngasilake bilangan G\"odel
saka~$!E(\gn{!E(x)})$. (Elinga, $\gn{!E(x)}$ iku numeral standar kanggo
bilangan G\"odel saka~$!E(x)$, yaiku $\num{\Gn{!E(x)}}$). Yen
$\Obj{diag}$ minangka simbol fungsi ing $\Th{T}$ kang merepresentasikake
fungsi $\fn{diag}$, kita bisa netepake $!A$ minangka formula
$!B(\Obj{diag}(\gn{!B(\Obj{diag}(x))}))$. Gatekna yen
\begin{align*}
\fn{diag}(\Gn{!B(\Obj{diag}(x))}) & = 
\Gn{!B(\Obj{diag}(\gn{!B(\Obj{diag}(x))}))} \\
& = \Gn{!A}.
\end{align*}
Yen $\Th{T}$ bisa !!{derive}
\[
\Obj{diag}(\gn{!B(\Obj{diag}(x))}) = \gn{!A},
\]
teori iku bisa !!{derive} $!B(\Obj{diag}(\gn{!B(\Obj{diag}(x))}))
\liff !B(\gn{!A})$. Nanging sisih kiwa, miturut
definisi, iku~$!A$.

Mesthi wae, $\Obj{diag}$ lumrahe dudu simbol fungsi saka
$\Th{T}$, lan mesthi dudu saka~$\Th{Q}$. Nanging, amarga $\fn{diag}$
komputabel, fungsi iki \emph{bisa direpresentasikake} ing~$\Th{Q}$
dening sawijining formula $!D_{\fn{diag}}(x,y)$. Mula tinimbang nulis
$!B(\Obj{diag}(x))$, kita bisa nulis
$\lexists[y][(!D_{\fn{diag}}(x,y) \land !B(y))]$. Saliyane panggantos
iki, rancangan bukti ing ndhuwur tetep lumaku, malah wis lumaku
ing~$\Th{Q}$.
\end{explain}

\begin{lem}
\ollabel{lem:fixed-point} Upamana $!B(x)$ formula sembarang kang mung
nduweni siji variabel bebas~$x$. Mula ana ukara~$!A$ saengga
$\Th{Q} \Proves !A \liff !B(\gn{!A})$.
\end{lem}


\begin{proof}
Kanggo $!B(x)$ kang diwenehake, tetepna $!E(x)$ minangka formula
$\lexists[y][(!D_{\fn{diag}}(x,y) \land !B(y))]$ lan tetepna $!A$
minangka diagonalisasine, yaiku formula $!E(\gn{!E(x)})$.

Amarga $!D_{\fn{diag}}$ merepresentasikake $\fn{diag}$, lan
$\fn{diag}(\Gn{!E(x)}) = \Gn{!A}$, $\Th{Q}$ bisa !!{derive}
\begin{align}
  & !D_{\fn{diag}}(\gn{!E(x)}, \gn{!A}) \ollabel{repdiag1} \\
  & \lforall[y][(!D_{\fn{diag}}(\gn{!E(x)},y) \lif
  \eq[y][\gn{!A}])]. \ollabel{repdiag2}
\end{align}
Saiki kita nuduhake yen $\Th{Q} \Proves !A \liff !B(\gn{!A})$. Kita
ngandharake kanthi informal, mung nggunakake logika lan fakta kang
!!{derivable} ing~$\Th{Q}$.

Sepisan, upamana~$!A$, yaiku $!E(\gn{!E(x)})$. Yen bali menyang
definisi $!E(x)$, kita weruh yen $!E(\gn{!E(x)})$ iku
\[
\lexists[y][(!D_{\fn{diag}}(\gn{!E(x)},y) \land !B(y))].
\]
Gatekna sawijining~$y$ kang kaya mangkono. Amarga
$!D_{\fn{diag}}(\gn{!E(x)},y)$, miturut \olref{repdiag2},
$y = \gn{!A}$. Mula saka $!B(y)$ kita entuk~$!B(\gn{!A})$.

Saiki upamana $!B(\gn{!A})$. Miturut \olref{repdiag1}, kita nduweni
\begin{align*}
& !D_{\fn{diag}}(\gn{!E(x)}, \gn{!A}) \land !B(\gn{!A}).
\intertext{Mula}
& \lexists[y][(!D_{\fn{diag}}(\gn{!E(x)},y) \land !B(y))].
\end{align*}
Nanging iki persis $!E(\gn{!E(x)})$, yaiku~$!A$.
\end{proof}

\begin{digress}
Bandhingna iki karo bukti lema titik tetep ing teori komputabilitas.
Bedane, ing kene kita arep nemtokake sawijining \emph{pratelan}
nganggo pratelan iku dhewe, dene ing kana kita arep nemtokake
sawijining \emph{fungsi} nganggo fungsi iku dhewe; sanajan ana
beda iki, gagasane sejatine padha.
\end{digress}

\begin{prob} 
  !!^a{formula}~$!A(x)$ iku \emph{definisi bebener} yen $\Th{Q}
  \Proves !B \liff !A(\gn{!B})$ kanggo saben !!{sentence}s~$!B$.
  Tuduhna nganggo lema titik tetep yen ora ana !!{formula} kang dadi
  definisi bebener.
\end{prob}

\end{document}
